\section{Certificates and the work used to produce them}
\label{sec:certificates}

Fix a compact axis-aligned box $D\subset\R^n$, a nonempty closed feasible
set $F\subset D$, and a continuous objective $f:D\to\R$. Write
\[
 f_* = \min_{y\in F}f(y),\qquad m(y)=f(y)-f_*,\qquad
 E(\eta)=\{y\in F:m(y)\leq\eta\}.
\]
The function $m$ is nonnegative on $F$; it can be negative outside $F$.
For each closed test box $B=\prod_i[l_i,u_i]\subset D$, the fixed
relaxation scheme supplies a set $R_B\supset F\cap B$ and a projected
objective $f_B$ with $f_B(y)\leq f(y)$ on $F\cap B$. Its bound is
\[
 L(B)=\inf_{z\in R_B}f_B(z),\qquad \inf\varnothing=+\infty.
\]
For a lifted relaxation, $f_B(y)$ is the infimum over auxiliary variables
at the fixed original point $y$. All statements concern exact certificates
for this specified scheme. They do not include numerical errors or the cost
of computing an individual bound.

\begin{definition}[Three spatial certificate classes]
\label{cert:def-classes}
A box is valid at additive tolerance $\epsilon>0$ when
$L(B)\geq f_*-\epsilon$. Let $N_{\mathrm{cover}}(\epsilon)$ be the minimum
size of a finite valid-box cover of $D$, and let
$N_{\mathrm{rect}}(\epsilon)$ be the minimum size of such a cover whose
box interiors are pairwise disjoint. Let $N_{\mathrm{tree}}(\epsilon)$ be
the minimum number of leaves of a finite axis-parallel split tree rooted
at $D$, all of whose leaf boxes are valid. These trees split a box into two
closed children at one coordinate value. A missing finite certificate
gives the corresponding minimum the value $+\infty$.
\end{definition}

Thus
\begin{equation}\label{cert:eq-benchmarks}
 N_{\mathrm{cover}}\leq N_{\mathrm{rect}}\leq N_{\mathrm{tree}}.
\end{equation}
An arbitrary rectangular partition need not itself be a split tree.
Refining it to a coordinate grid may increase its size. We will obtain
lower bounds for covers and upper bounds by constructing trees, which can
give the same asymptotic rate for these three different finite minima.

A complete binary tree with $\ell$ leaves has $2\ell-1$ nodes. A complete
$b$-ary tree has $1+b(\ell-1)/(b-1)$ nodes. These identities count tree
nodes only. They count neither repeated relaxation evaluations nor work
performed by tightening or probing. A run with a feasible incumbent
$U\geq f_*$ may prune at $L(B)\geq U-\epsilon$; every such pruning test is
valid in \cref{cert:def-classes}. Upper bounds below use the ideal incumbent
$U=f_*$ and therefore measure certification after an optimal solution is
available. They do not measure the work needed to find that solution.

\subsection{Quadratic gaps and owned tests}

For $y\in B$, put
\begin{equation}\label{cert:eq-q}
 a_i^B(y)=(y_i-l_i)(u_i-y_i),\qquad
 q_B(y)=\sum_{i=1}^n a_i^B(y),\qquad
 d_i^B(y)=\min\{y_i-l_i,u_i-y_i\}.
\end{equation}
We will use $a_i^B\geq(d_i^B)^2$ and
$a_i^B\leq(u_i-l_i)d_i^B$. If $C\subset B$, then $q_C\leq q_B$ on $C$.
Zero-width factors are zero. Dimension-collapsing reductions are permitted:
the surviving degenerate box remains a node, evaluated in its affine hull.

\begin{assumption}[Objective lower gaps]\label{cert:ass-gap}
For a fixed $\alpha>0$, the bound gap is
\begin{equation}\label{cert:eq-bound-gap}
 L(B)\leq f(y)-\alpha q_B(y)\qquad(y\in F\cap B).
\end{equation}
The stronger pointwise gap is
\begin{equation}\label{cert:eq-point-gap}
 f_B(y)\leq f(y)-\alpha q_B(y)\qquad(y\in F\cap B).
\end{equation}
Both are required uniformly over the boxes allowed by the model.
\end{assumption}

Every valid raw test satisfies
\begin{equation}\label{cert:eq-owned-valid}
 m(y)+\epsilon\geq\alpha q_C(y)\qquad(y\in F\cap C).
\end{equation}
After tightening, the same inequality may hold only on a subset of the
closed test box. To handle this distinction, an \emph{owned test} is a
pair $(C,A)$, where $C$ is a closed box and $A\subset C$ is a rectangle
with possibly open endpoints. Its validity means
\cref{cert:eq-owned-valid} for $y\in F\cap A$. Owners are Borel and convex.
An owned certificate is a finite such family whose owners cover $F$.
It is a different certificate class from \cref{cert:def-classes}.

\begin{lemma}[Boundary-safe event accounting]\label{cert:lem-events}
Consider a finite terminating run using axis-parallel splits and the
following reductions. An original-feasibility reduction $B\mapsto B'$
retains $F\cap B$. A same-relaxation reduction retains
\[
 \{z\in R_B:f_B(z)\leq U-\epsilon\},\qquad U\geq f_*.
\]
An emptying reduction is a terminal event. Expand any stopping bound
obtained as a minimum of probe-child bounds into the finite child
partition on which those bounds rely, recursively until raw successful
tests are reached. Let $\ell_{\mathrm{aug}}$ count the terminal events in
this augmented event tree, and let $R_{\mathrm{rel}}$ count nonempty
same-relaxation reduction rounds. Under the pointwise gap there is an
owned certificate for $F$ with
\begin{equation}\label{cert:eq-event-count}
 K_F\leq\ell_{\mathrm{aug}}+2nR_{\mathrm{rel}}.
\end{equation}
The bound gap alone suffices when there are no same-relaxation reductions.
Original-feasibility reductions add no owners meeting $F$.
\end{lemma}

\begin{proof}
Initially the root owns $D$. At a split in coordinate $i$ at $s$, assign
$A\cap\{y_i\leq s\}$ to the lower child and
$A\cap\{y_i>s\}$ to the upper child. Suppose next that the current
closed test is $B=\prod_i[l_i,u_i]$ and a nonempty reduction retains
$B'=\prod_i[l'_i,u'_i]$. Carry $A\cap B'$ forward, including its boundary.
For each $i$, assign the discarded lower and upper owners
\[
 A_i^- = A\cap\bigcap_{j<i}\{l'_j\leq y_j\leq u'_j\}
             \cap\{y_i<l'_i\},\qquad
 A_i^+ = A\cap\bigcap_{j<i}\{l'_j\leq y_j\leq u'_j\}
             \cap\{y_i>u'_i\}.
\]
Use the closed slabs
\[
 C_i^- =\prod_{j<i}[l'_j,u'_j]\times[l_i,l'_i]
                  \times\prod_{j>i}[l_j,u_j]
\]
and their upper analogues as tests. Empty owners are omitted. The first
coordinate outside $B'$ gives a disjoint partition of $A\setminus B'$,
with at most $2n$ convex rectangular owners. This remains true when
$B'$ has zero widths. An emptying reduction uses $(B,A)$ as one terminal
event. Induction shows that all terminal owners partition $D$.

A discarded feasibility owner contains no feasible point. If a feasible
$y$ is discarded by a same-relaxation reduction, then $y\in R_B$ and
$f_B(y)>U-\epsilon\geq f_*-\epsilon$. For its slab $C\subset B$,
\[
 f(y)-\alpha q_C(y)\geq f(y)-\alpha q_B(y)
                         \geq f_B(y)>f_*-\epsilon.
\]
At a successful raw terminal bound from a containing source box $B$,
the bound gap similarly gives
$f(y)-\alpha q_C(y)\geq L(B)\geq f_*-\epsilon$.
A relaxed-infeasibility test owns no feasible point. The expansion of
probing bounds ensures that each terminal successful test has such a raw
source. Filtering the partition by nonempty intersection with $F$ proves
the assertion and the count.
\end{proof}

Closed discarded slabs need not be valid at their retained boundary.
For example, let $D=[-1,1]\times[0,1]$, $F=\{0\}\times[0,1]$, $f=0$,
$R_B=F\cap B$, and $f_B=f-\alpha q_B$ on $R_B$. Exact feasibility
reduction retains $\{0\}\times[0,1]$. Both closed frame slabs contain
$(0,1/2)$, where their $q$ is $1/4$, although their discarded owners
contain no feasible point. This example satisfies the pointwise gap.
It shows why \cref{cert:lem-events} gives an owned certificate and does
not imply $N_{\mathrm{cover}}\leq K_F$.

\subsection{Which costs the accounting controls}

Let $T_{\mathrm{aug}}$ count processed nodes in the augmented event tree.
Then \cref{cert:eq-event-count} gives
\begin{equation}\label{cert:eq-work}
 K_F\leq T_{\mathrm{aug}}+2nR_{\mathrm{rel}}.
\end{equation}
Consequently a geometric lower bound $K_F\geq H(\epsilon)$ bounds this
combined event and reduction work. If at most $r$ relevant reduction
rounds occur per augmented node, it also gives
$T_{\mathrm{aug}}\geq H(\epsilon)/(1+2nr)$. Without that hypothesis it
does not give a node lower bound. When no stopping test uses a probe
partition, actual search nodes provide $T_{\mathrm{aug}}$. Otherwise
virtual probe tests must be included or separately charged: a root can
certify itself using many child probes while committing no split.

There is a separate evaluation-cost conclusion. Suppose every relevant
reduction round is charged to a distinct counted relaxation evaluation,
and every raw source needed by a successful terminal bound, including
probe sources, is counted among $S_{\mathrm{eval}}$ evaluations. Keep
each successful source box once, and keep the discarded slab owners.
This new family covers $F$ by valid source boxes and valid owned subsets,
possibly with overlap. Its count satisfies
\begin{equation}\label{cert:eq-eval}
 K_{\mathrm{eval}}\leq S_{\mathrm{success}}+2nR_{\mathrm{rel}}
                         \leq(2n+1)S_{\mathrm{eval}}.
\end{equation}
Every surviving feasible point is covered by a successful source; every
discarded feasible point has a slab owner. A source box satisfies
\cref{cert:eq-owned-valid} everywhere on its feasible subset, so the
per-test lower bounds below apply to this new family. The count
$K_{\mathrm{eval}}$ is separate from $K_F$: deduplicating inherited
evaluations need not decrease the number of terminal owners. Uncharged
reductions or uncharged probe bounds are outside
\cref{cert:eq-eval}. Tightening rounds, removed-piece certificates, and
relaxation evaluations remain distinct quantities.
